% Options for packages loaded elsewhere
\PassOptionsToPackage{unicode}{hyperref}
\PassOptionsToPackage{hyphens}{url}
\PassOptionsToPackage{dvipsnames,svgnames,x11names}{xcolor}
\documentclass[
  11pt,
  a4paper,
]{article}
\usepackage{xcolor}
\usepackage[margin=24mm]{geometry}
\usepackage{amsmath,amssymb}
\setcounter{secnumdepth}{-\maxdimen} % remove section numbering
\usepackage{iftex}
\ifPDFTeX
  \usepackage[T1]{fontenc}
  \usepackage[utf8]{inputenc}
  \usepackage{textcomp} % provide euro and other symbols
\else % if luatex or xetex
  \usepackage{unicode-math} % this also loads fontspec
  \defaultfontfeatures{Scale=MatchLowercase}
  \defaultfontfeatures[\rmfamily]{Ligatures=TeX,Scale=1}
\fi
\usepackage{lmodern}
\ifPDFTeX\else
  % xetex/luatex font selection
  \setmainfont[]{Cambria}
  \setmathfont[]{Cambria Math}
\fi
% Use upquote if available, for straight quotes in verbatim environments
\IfFileExists{upquote.sty}{\usepackage{upquote}}{}
\IfFileExists{microtype.sty}{% use microtype if available
  \usepackage[]{microtype}
  \UseMicrotypeSet[protrusion]{basicmath} % disable protrusion for tt fonts
}{}
\makeatletter
\@ifundefined{KOMAClassName}{% if non-KOMA class
  \IfFileExists{parskip.sty}{%
    \usepackage{parskip}
  }{% else
    \setlength{\parindent}{0pt}
    \setlength{\parskip}{6pt plus 2pt minus 1pt}}
}{% if KOMA class
  \KOMAoptions{parskip=half}}
\makeatother
\setlength{\emergencystretch}{3em} % prevent overfull lines
\providecommand{\tightlist}{%
  \setlength{\itemsep}{0pt}\setlength{\parskip}{0pt}}
\pagestyle{plain}
\setlength{\emergencystretch}{3em}
\setlength{\parskip}{0.35em}
\IfFileExists{bookmark.sty}{\usepackage{bookmark}}{\usepackage{hyperref}}
\IfFileExists{xurl.sty}{\usepackage{xurl}}{} % add URL line breaks if available
\urlstyle{same}
\hypersetup{
  pdftitle={Correction to Hormesis as a Geometric Necessity of Bounded Adaptive Systems},
  pdfauthor={Daniel John Murray},
  colorlinks=true,
  linkcolor={black},
  filecolor={Maroon},
  citecolor={Blue},
  urlcolor={blue},
  pdfcreator={LaTeX via pandoc}}

\title{Correction to Hormesis as a Geometric Necessity of Bounded Adaptive Systems}
\usepackage{etoolbox}
\makeatletter
\providecommand{\subtitle}[1]{% add subtitle to \maketitle
  \apptocmd{\@title}{\par {\large #1 \par}}{}{}
}
\makeatother
\subtitle{Proposed author corrigendum for DOI 10.1177/15593258261469171}
\author{Daniel John Murray}
\date{12 October 2026}

\begin{document}
\maketitle

Independent Researcher, Melbourne, Australia. ORCID: 0009-0005-1794-5945.

This proposed correction concerns Murray DJ, \emph{Dose Response}. 2026;24(3):15593258261469171, \href{https://doi.org/10.1177/15593258261469171}{\nolinkurl{doi:10.1177/15593258261469171}}. It corrects the mathematical necessity, peak-location, finite-change attenuation, approximation and zone-width statements. It does not constitute an accepted or issued journal correction.

\textbf{Necessity and baseline normalization.} In Sections 3.3, 7.2, 7.5 and the Conclusion, activation before toxicity and greater high-dose damage are insufficient to guarantee a positive-dose beneficial interval. The adopted Möbius composition law supplies an additive rapidity coordinate; it does not choose the relative amplitude or slope of the two dose-dependent increments. Section 2.2 correctly states that Aczél's general theorem alone does not select that coordinate, and this clarification is retained.

For the published primary model with equal linear-dose logistic slope \(s\), repair amplitude \(\alpha\), damage amplitude \(\beta>\alpha\), and \(0<D_a<D_t\), set \(a=e^{sD_a}\) and \(t=e^{sD_t}\). A beneficial interval relative to the actual zero-dose control exists if and only if

\[\frac{\alpha a}{(1+a)^2}>\frac{\beta t}{(1+t)^2}.\]

Zero-dose logistic tails should be subtracted from the rapidity increments, or the response should be divided by its actual zero-dose value, so the control is exactly one. The dose-dependent peak is unchanged by this normalization. With \(s=1,D_a=4,D_t=5,\alpha=0.3,\beta=1\), the stated threshold ordering and high-dose inequality hold, but the response decreases throughout the positive-dose domain. The geometric-necessity claim and the corresponding mathematical exclusion of a monotonically harmful endpoint must therefore be withdrawn at their stated generality.

\textbf{Peak location.} Section 5.1, Prediction 1, Section 7.6, the Conclusion and Appendix A.2 should not state a universal peak near \(2D_a\). For the equal-slope primary model, an interior peak exists under the condition above and is exactly

\[D_*=\frac1s\log\left[
\frac{\sqrt{\alpha a}\,t-\sqrt{\beta t}\,a}
{\sqrt{\beta t}-\sqrt{\alpha a}}
\right].\]

The peak depends on both amplitudes, both thresholds and the slope. Unequal slopes require their own stationary equation. The later multiplicative Hill response has a different exact law and is not an algebraic replacement for the published linear-dose rapidity model.

\textbf{Attenuation and finite changes.} Sections 4.2--4.3, Predictions 3 and 5, related amplitude conclusions and Appendix A.3 extend a local metabolic-control approximation to finite changes. The summation theorem constrains the sum of local logarithmic coefficients. It neither forces equal control coefficients nor determines a finite doubling. In the explicitly chosen equal-resistance series model, increasing one rate by a fraction \(A_{\rm mol}\) gives

\[A_{\rm func}=\frac{A_{\rm mol}}{n+(n-1)A_{\rm mol}},\qquad
\eta=n+(n-1)A_{\rm mol}.\]

The approximation \(\eta\simeq n\) is its small-change limit. Doubling one rate gives a flux increase of \(1/3\) for two equal steps and \(1/5\) for three, rather than \(1/2\) and \(1/3\). Other kinetic models give other finite responses, so step counts alone do not fix the observed attenuation.

The published comparisons with broad hormesis features and activation--toxicity ordering do not validate the withdrawn universal peak or finite-amplitude statements. The bounded-composition framework and conditional response models remain available. The companion mathematical amendment provides complete proofs, exact model-specific onset, peak and return-to-control laws, and independent reproducibility checks. The journal's handling and final wording of a formal corrigendum remain subject to editorial review.

\textbf{Coincident thresholds and locally linear harm.} Section 7.2's general near-coincident-threshold guarantee is withdrawn. Identical repair and damage sigmoids with equal midpoint and greater damage give no beneficial interval. In the common-slope logistic model at fixed amplitudes, sufficiently small threshold gaps also give none. Different sigmoid shapes require their own amplitude--slope conditions. A negative linear low-dose term is permitted in the common-slope logistic model and in the additive Hill model at exponent one. The composition law therefore does not exclude locally linear harm; this does not identify a globally proportional dose law.

\textbf{Approximation bounds and zone width.} Section 3.4 and Appendix A.1 require a full calibration and error definition for the product shortcut. Repair-only saturation matching and tanh concavity do not establish a conservative bound for the combined repair--damage response. The full product can overestimate the exact response even with a positive actual baseline. The uniform conservative interpretation and the unqualified error bound are withdrawn. The peak and zone-width ranges in Section 5.2, Table 2 and Appendix A.2 also depend on an incompletely stated steepness and control protocol. One compatible reconstruction, with each logistic slope set to \(4/D_j\), reproduces Appendix A.2's range 1.69--2.43; this does not identify the original unstated protocol uniquely, and other slopes give other peaks. A detectable zone must use the actual control, a declared positive threshold and both threshold crossings. No universal 8--80-fold width follows from bounded composition.

\textbf{Finite amplitudes and control.} This also amends the peak-range, mean and aggregate-agreement claims in the abstract, Section 5.3 equation (8), Section 6.4 and Table 2. For the stated molecular changes \(A_{\rm mol}\in[0.8,1.5]\) and \(n\in\{2,3\}\), the corrected series-model flux factors span 117.4--142.9\% of control. These are conditional finite flux ceilings, not a predicted distribution of hormesis peaks. A normalized bounded viability with actual control \(p_0\) is below \(1/p_0\) at finite rapidity, so its amplitude also requires an independently measured baseline and a flux-to-endpoint bridge.

\section{References}\label{references}

\begin{enumerate}
\def\labelenumi{\arabic{enumi}.}
\tightlist
\item
  Murray DJ. \emph{Dose Response}. 2026;24(3):15593258261469171. \href{https://doi.org/10.1177/15593258261469171}{\nolinkurl{doi:10.1177/15593258261469171}}.
\item
  Kacser H, Burns JA. \emph{Symp Soc Exp Biol}. 1973;27:65-104.
\item
  Heinrich R, Rapoport TA. \emph{Eur J Biochem}. 1974;42:89-95. \href{https://doi.org/10.1111/j.1432-1033.1974.tb03318.x}{\nolinkurl{doi:10.1111/j.1432-1033.1974.tb03318.x}}.
\end{enumerate}

\end{document}
